% =====================================================================
%  CV TEMPLATE — Nishal Stanislaus Silva, PhD
%  Fill from knowledge-graph.md (job-search project). Placeholders look like
%      <<KG §x.y | what to put here | rule>>
%  Compile: pdflatex (or xelatex). Derived from the five CV variants in the
%  job-search project (generic, ai-industry, ai-research, musictech, long).
% =====================================================================
%
%  HOW TO USE
%  1. Read the JD. Pick ONE variant from the VARIANT PROFILES table below.
%  2. That row fixes: tagline, section order, competency group names, stack
%     lines, publication count, and page target. Do not mix rows.
%  3. Fill every <<KG ...>> placeholder from the knowledge graph.
%  4. Reorder/delete sections per the row's "Order" column (each section below
%     is a self-contained block between ==== rules; cut and paste whole blocks).
%  5. Output a changelog: variant chosen, why, what was cut, final page count.
%
% =====================================================================
%  VARIANT PROFILES  (choose exactly one)
% =====================================================================
%
%  V1  AI-INDUSTRY  — applied/product ML at a company. DEFAULT for most roles.
%      Tagline : Machine Learning Research Engineer | Deep Learning, Time-Series,
%                Pattern Detection | Applied ML
%      Order   : Summary, Competencies, Stack, Education, Portfolio(4 items),
%                Experience, Publications(3), Highlights
%      Groups  : ML Engineering / Domain Expertise / Research & Leadership
%      Stack   : Languages / ML & AI / Data & Cloud / Dev & MLOps
%                (emphasise Kubernetes, MLflow, REST APIs, ONNX, SageMaker,
%                 LLMs, RAG, recommender & ranking systems)
%      Pages   : 2      Cut: Datasets, Peer Review, Projects, Talks, Supervision
%
%  V2  AI-RESEARCH  — research scientist / applied scientist / national lab.
%      Tagline : Research Scientist | Real-Time Multimodal ML, Benchmarks &
%                Open Datasets
%      Order   : Summary, Competencies, Stack, Education, Experience,
%                Publications(FULL), Datasets(FULL), Peer Review, Talks,
%                Supervision, Highlights
%      Groups  : ML Engineering / Domain Expertise / Research & Leadership
%      Stack   : as V1 but lead ML & AI with JAX, HuggingFace, PyTorch
%      Pages   : 3      Cut: Portfolio (links live in Publications instead)
%
%  V3  EDGE / EMBEDDED  — on-device ML, TinyML, real-time systems, silicon.
%      Tagline : ML Research Engineer | Real-Time & Edge Inference | Audio,
%                Sensor & Time-Series
%      Order   : Summary, Competencies, Stack, Portfolio, Education, Experience,
%                Datasets(short), Publications(5-7), Projects, Highlights
%      Groups  : ML Engineering / Edge & System Optimization / Research Engineering
%      Stack   : Languages / ML & AI / EDGE & SYSTEMS (Raspberry Pi, Embedded
%                Linux, JUCE, Elk Audio OS, ARM profiling, OpenCV) / Data & Cloud
%                / Dev & MLOps
%      Pages   : 2, or 3 if the JD asks for publications
%
%  V4  MUSIC-TECH  — audio ML, MIR, instruments, creative tools.
%      Tagline : Music Technology ML Research Engineer | Audio, DSP &
%                Real-Time Systems
%      Order   : Summary, PORTFOLIO & DEMOS (high, right after summary),
%                Competencies, Stack, Experience, Publications(4, include the
%                PhD thesis), EDUCATION (moved near the end), Highlights
%      Groups  : Music Tech & Audio ML / ML Engineering / Research &
%                Cross-Domain Depth
%      Stack   : Languages / ML & AI / AUDIO & SIGNALS (Librosa, SciPy, MFCC,
%                STFT, JUCE, Elk Audio OS, MIDI) / Data & Cloud / MLOps & Dev
%      Pages   : 2. Keep the guitarist line in Highlights (it is an asset here).
%      Note    : the old musictech CV dropped the PR/sponsorship line — keep it.
%
%  V5  COMPUTER VISION / INDUSTRIAL  — inspection, manufacturing, robotics, retail.
%      Tagline : Computer Vision & ML Engineer | Industrial Inspection, IoT &
%                Edge Deployment
%      Order   : Summary, Competencies, Stack, Education, Experience (MAS gets
%                4-5 bullets, postdoc trimmed to 2), Projects(CV-heavy),
%                Publications(3), Highlights
%      Groups  : Computer Vision & Perception / ML Engineering / Systems &
%                Deployment
%      Stack   : swap the audio line for COMPUTER VISION (OpenCV, template &
%                feature matching, OCR, image registration, industrial cameras,
%                Cognex VisionPro, HALCON)
%      Pages   : 2      Lead projects with P14/P15/P16/P18/P19 (KG §8.2).
%
%  V6  GENERIC  — speculative applications, recruiters, unclear JDs.
%      Tagline : Research Engineer | Real-Time ML, DSP & Time-Series |
%                Music Technology
%      Order   : Summary, Competencies (FLAT comma list, no group headings),
%                Education, Stack, Experience, Publications(3), Highlights
%      Pages   : 2
%
% =====================================================================
%  RULES THAT APPLY TO EVERY VARIANT
% =====================================================================
%  * Never write ">90% F1". Cite F1 = 0.76 (KG §16).
%  * MAS Holdings is ONE entry, Jan 2015 - Jul 2020 (KG §16).
%  * Promptly: alignment algorithm, RIP integration, mechanism consulting.
%    Never mention the patent - Nishal is not a listed inventor (KG §8.2).
%  * Verified numbers before self-reported ones (KG §13).
%  * Languages: English native; add Sinhala/Italian/French only per KG §15.6.
%  * Keep the PR + no-sponsorship line for all Canadian roles.
%  * Never exceed 3 pages. If it overflows: cut Talks, then Projects, then
%    Peer Review - cut Publications last.
% =====================================================================

\documentclass[11pt]{article}

% --------- PACKAGES ----------
\usepackage[left=0.7in,top=0.7in,right=0.7in,bottom=0.7in]{geometry}
\usepackage{enumitem}
\usepackage[hidelinks]{hyperref}
\usepackage{titlesec}
\usepackage{array}
\usepackage{setspace}
\usepackage[T1]{fontenc}
\usepackage{newtxtext,newtxmath}   % Times-like
\input{glyphtounicode}
\pdfgentounicode=1

\setstretch{1.05}
\pagenumbering{gobble}

% --------- SECTION FORMAT ----------
\titleformat{\section}{\Large\scshape}{}{4em}{}[\vspace{-0.7em}\rule{\linewidth}{0.4pt}\vspace{-0.4em}]
\titlespacing*{\section}{0pt}{1em}{0.4em}

% \begin{cvrole}{Title}{Organisation}{City, Country}{Start}{End}
\newenvironment{cvrole}[5]{%
  \par\noindent\textbf{\large #1} | \textit{\small #2, #3}\hfill \textit{\small #4 -- #5}\par
  \begin{itemize}[leftmargin=2em, itemsep=-0.3em, topsep=-0.5em]%
}{%
  \end{itemize}\par\noindent\mbox{}\par\vspace{0.1em}%
}

% --------- FIXED IDENTITY (KG §1 — do not edit) ----------
\newcommand{\cvName}{Nishal Stanislaus Silva}
\newcommand{\cvEmail}{nishal.silva@hotmail.com}
\newcommand{\cvPhone}{+1 613 200 2030}
\newcommand{\cvCity}{Toronto, ON}
\newcommand{\cvSite}{https://nishal.xyz}
\newcommand{\cvLinkedIn}{https://www.linkedin.com/in/nishal-silva}
\newcommand{\cvGitHub}{https://github.com/ns2max}
\newcommand{\cvStatus}{Canadian Permanent Resident, Eligible to work in Canada without sponsorship}

% --------- PER-APPLICATION ----------
% Paste the chosen variant's tagline verbatim; use the SAME one in the cover letter.
\newcommand{\cvTagline}{<<VARIANT | tagline from the chosen profile row>>}

% =====================================================================
\begin{document}
\pagestyle{empty}

% ============================== HEADER ===============================
% [ALL VARIANTS] Never reorder or trim. Drop \cvGitHub only if space is tight
% and the JD has no code component.
\begin{center}
    {\LARGE \scshape \textbf{\cvName}}{\large \textit{, PhD}}
    \\[1pt]
    {\small \cvTagline}
    \\[1pt]
    {\footnotesize\textit{\cvStatus}}
    \\[1pt]
    {\small
    \href{mailto:\cvEmail}{\cvEmail}
    \,\,|\,\, \cvPhone
    \,\,|\,\, \cvCity
    \,\,|\,\, \href{\cvSite}{nishal.xyz}
    \,\,|\,\, \href{\cvLinkedIn}{linkedin.com/in/nishal-silva}
    \,\,|\,\, \href{\cvGitHub}{github.com/ns2max}}
\end{center}

% ========================= PROFESSIONAL SUMMARY ======================
% [ALL VARIANTS] Position 1 always. 4-5 sentences.
\section*{Professional Summary}

<<KG §0 + §13 | Sentence 1: archetype + 10+ years across industry and academia.
Sentence 2: what he builds end-to-end, in the JD's vocabulary (KG §14).
Sentence 3: the domain spread that matches the JD (V1 real-time/streaming ML,
time-series, CV, structured data | V3 edge inference and system-level efficiency
| V4 MIR, smart instruments, multimodal audio+sensor | V5 industrial inspection,
IoT, retail analytics).
Sentence 4: verified proof (14 ms on RPi4, F1 0.76, 74x faster than DTW) plus one
industrial figure (up to 300\% efficiency, 99.5\% inspection-time reduction).
Sentence 5: publication credibility (JAES, IEEE, ACM) + the JD-specific hook.
| never ">90\% F1">>

% ========================== CORE COMPETENCIES ========================
% [ALL VARIANTS] Group headings come from the chosen variant row.
% V6 GENERIC: delete the three \textbf lines and use ONE flat comma-separated
% paragraph instead (no group headings) — that is what cv-generic.md does.
\section*{Core Competencies}

\textbf{<<VARIANT | group 1 heading>>:} <<KG §5.2, §5.7 | 8-12 items, JD-ordered>>\\
\textbf{<<VARIANT | group 2 heading>>:} <<KG §5.3-5.6 | 8-12 items>>\\
\textbf{<<VARIANT | group 3 heading>>:} <<KG §5.7 + §13 | 6-10 items>>

% ============================ TECHNICAL STACK ========================
% [ALL VARIANTS] Line 3 is the variant-specific line — see the profile row.
% Only list KG depth-W items when the JD names them explicitly (KG §5).
\section*{Technical Stack}

\textbf{Languages:} <<KG §5.1 | Python, C++, C, MATLAB, JavaScript, Shell, ...>>\\
\textbf{ML / AI:} <<KG §5.2 | V1/V5: TensorFlow, Keras, PyTorch, scikit-learn, ONNX, SageMaker | V2: JAX, HuggingFace, PyTorch first | V3: add model compression, quantisation | V4: keep short, audio line carries the weight>>\\
\textbf{<<VARIANT | V1/V2 "Signal \& Audio Processing" | V3 "Edge \& Systems" | V4 "Audio \& Signals" | V5 "Computer Vision">>:} <<KG §5.3 / §5.4 / §5.5 | pick the matching subsection>>\\
\textbf{Data \& Cloud:} <<KG §5.6 | AWS (EC2, S3, ECS, MWAA, RDS), Snowflake, SQL, ETL, NumPy, Pandas>>\\
\textbf{Dev \& MLOps:} <<KG §5.6 | Docker, Airflow, Jenkins, Git, Linux/UNIX, pytest, CI/CD | V1 add Kubernetes, MLflow, REST APIs>>

% =============================== EDUCATION ===========================
% [ALL VARIANTS] Position: V1/V2/V5 after Stack | V6 BEFORE Stack |
% V4 move this whole block near the end, after Publications.
\section*{Education}
\textbf{Ph.D -- Information Engineering and Computer Science} \textit{\small{ -- University of Trento, Italy}} \hfill \textit{Jan 2025}\\[2pt]
\noindent\textbf{M.Sc -- Telecommunication and Electronic Engineering} \textit{\small{ -- Sheffield Hallam University, UK}} \hfill \textit{Aug 2018}\\[2pt]
\noindent\textbf{B.Eng (Hons) -- Electronic Engineering} \textit{\small{ -- Sheffield Hallam University, UK}} \hfill \textit{Mar 2014}
% V4 MUSIC-TECH may append:  \\[2pt]
% \noindent\textbf{Rock \& Pop Guitar, Grade 8} \textit{\small{ -- Trinity College London, UK}} \hfill \textit{Jul 2015}

% ========================== PORTFOLIO AND DEMOS ======================
% [V1 keep 3-4 items] [V3 keep, lead with Nebula] [V4 MOVE THIS BLOCK UP,
% directly after the Summary, and keep all items] [V2/V5/V6 DELETE].
\section*{Portfolio and Demos \footnotesize {\normalfont{(full portfolio: \href{https://nishal.xyz/research}{\textit{nishal.xyz/research}})}}}
\begin{itemize}[leftmargin=1.2em, itemsep=-0.25em]
    \item An Internet of Musical Things Performance Ecosystem: \href{https://youtu.be/jfm5L3DYIcc}{\textit{youtu.be/jfm5L3DYIcc}}
    \item Multisensory concert enabled by smart musical instruments: \href{https://youtu.be/axEHkdnxFB8}{\textit{youtu.be/axEHkdnxFB8}}
    \item Hot Licks: real-time pattern detection for MIDI instruments (MIDI Innovation Awards 2023 finalist): \href{https://youtu.be/gkOqfOT8zXM}{\textit{youtu.be/gkOqfOT8zXM}}
    \item Nebula: C++17 audio feature extraction library with ARM latency benchmarks: \href{https://github.com/ns2max/nebula}{\textit{github.com/ns2max/nebula}}
    \item LiveLaTeX: real-time LaTeX compiler/preview extension for VS Code: \href{https://marketplace.visualstudio.com/items?itemName=ns2max.livelatex}{\textit{marketplace.visualstudio.com/items?itemName=ns2max.livelatex}}
\end{itemize}

\newpage
% ======================= PROFESSIONAL EXPERIENCE =====================
% [ALL VARIANTS] Bullet budget: 2 pages -> 3 bullets max per role;
% 3 pages -> up to 4. V5 inverts the weighting: MAS gets 4-5 bullets and the
% postdoc drops to 2. Section title: "Professional Experience" (V1/V5) or
% "Research and Work Experience" (V2/V3/V4/V6).
\section*{Professional Experience}

% ---- OPTIONAL gap cover (KG §4.6). Include for V1/V3/V5; for V2 the papers
% ---- in press already cover the period, so it is usually unnecessary.
% \begin{cvrole}{Independent Research \& Open-Source}{Self-directed}{Toronto, Canada}{Jan 2026}{Present}
%     \item <<KG §4.6 | Nebula: C++17 audio feature library with per-feature ARM latency benchmarking>>
%     \item <<KG §4.6 | LiveLaTeX: published VS Code extension (Tectonic-based live preview)>>
%     \item <<KG §4.6 | MIRaaS paper accepted at IEEE IS2 2026; Smart Drums in press at JAES>>
% \end{cvrole}

\begin{cvrole}{Postdoctoral Researcher}{University of Trento}{Trento, Italy}{Jan 2025}{Dec 2025}
    \item <<KG §4.1 | end-to-end pipeline: streaming audio + sensor + gesture data, acquisition -> DSP/feature engineering -> training/evaluation -> low-latency edge deployment with monitoring; EU-funded MUSMET (EIC Pathfinder)>>
    \item <<KG §4.1 + §13 | metric bullet: real-time inference at 14 ms on Raspberry Pi 4, F1 0.76, 74x faster than the DTW baseline; 31.4\% CPU within an embedded budget (JAES 2026)>>
    \item <<KG §4.1 | pick ONE by variant: V1/V2 benchmark suites (baselines + ablations) quantifying architecture/feature trade-offs | V2/V3 multimodal sync pipeline: audio + BCI/EEG + MR head/hand tracking synchronised across 4 musicians, with SAE Barcelona and g.tec | V3 MIRaaS edge-to-server offload and latency characterisation | V4 multisensory concerts and the Hot Licks Mapper routing tool>>
    \item <<KG §4.1 | 3-page only: partner coordination and 4 peer-reviewed outputs in 2025; or the teaching-tool / accessibility research thread (KG P3c) for ed-tech and accessibility JDs>>
\end{cvrole}

\begin{cvrole}{Visiting Researcher}{McGill University}{Montreal, Canada}{May 2024}{Aug 2024}
    \item <<KG §4.2 | multimodal IMU + audio gesture recognition fused by a hierarchical state machine; F1 0.78 across a 500-event corpus (IEEE IS2 2025)>>
    \item <<KG §4.2 | compute and power profiling for edge deployment feasibility; VAE- and diffusion-based synthetic data for limited-label regimes>>
\end{cvrole}

\begin{cvrole}{Visiting Researcher}{University of Visual and Performing Arts}{Colombo 02, Sri Lanka}{Jan 2024}{Mar 2024}
    \item <<KG §4.3 | human-centred evaluation of ML-driven interactive systems, combining performance metrics with user feedback to produce usability, interpretability and adoption recommendations (IEEE I3DA 2025)>>
\end{cvrole}

\begin{cvrole}{Technical Consultant}{Forestpin (Pvt) Ltd}{Colombo 05, Sri Lanka}{Aug 2020}{Feb 2021}
    \item <<KG §4.4 | ML and statistical modelling pipelines for financial forensics, compliance monitoring and risk detection on large structured enterprise datasets>>
    \item <<KG §4.4 | SQL + Python backend scoring services for anomaly flagging and production triage; translated compliance requirements into technical specifications>>
\end{cvrole}

\begin{cvrole}{Research Engineer}{MAS Holdings (Pvt) Ltd}{Colombo, Sri Lanka}{Jan 2015}{Jul 2020}
    \item <<KG §4.5 | end-to-end ML and computer-vision systems for manufacturing at scale; camera, sensor and IoT streams integrated into production workflows; operational efficiency improved by up to 300\%>>
    \item <<KG §4.5 | pick by variant: V5 care-label QC (OpenCV template + feature matching, conveyor with PLC reject, inspection time cut 99.5\%) | V5 loom-side fabric inspection (microscopic camera array, operator headcount cut 50\%) | V1/V5 Promptly on-demand printing (geometry-alignment algorithm, RIP integration, mechanism consulting -- no patent claim) | V5 lace archive with ERP-integrated search, still in production use>>
    \item <<KG §4.5 | real-time ingestion and ETL for high-frequency IoT and operational time-series across distributed sites in South Asia, North America and Africa; C++/Python inference engines on embedded hardware>>
    \item <<KG §4.5 | 3-page or V5 only: human-in-the-loop inspection models balancing throughput against risk; mentored engineers; established CI/CD and code-review practice; delivered group-wide training>>
\end{cvrole}

% ========================= SELECTED PUBLICATIONS =====================
% Count by variant: V1/V5/V6 -> 3   V4 -> 4 (include the PhD thesis)
%                   V3 -> 5-7        V2 -> FULL list, retitle "Publications"
%                                        and drop the "full list" footnote.
% V2 also appends the two theses:
%   Silva, N. Embedded Real-time Musical Pattern Detection for Smart Musical
%   Instruments. PhD Thesis, University of Trento, Jan. 2025.
%   Silva, N. On Musical Onset Detection via the S-Transform. MSc Thesis,
%   Sheffield Hallam University, 2017.
% Ordering: keep reverse-chronological. Delete the entries you do not use.
\section*{Selected Publications \footnotesize {\normalfont{(full list: \href{https://nishal.xyz/publications}{\textit{nishal.xyz/publications}})}}}
\begin{itemize}[leftmargin=1.2em, itemsep=-0.25em]
    \item Silva, N. and Turchet, L. \textit{Real-Time Audio Pattern Detection for Smart Musical Instruments}. Journal of the Audio Engineering Society, Vol.~74, No.~3, Mar. 2026. \textit{(F1 0.76 at 14 ms on Raspberry Pi 4, 74$\times$ faster than baseline)}
    \item Silva, N. and Turchet, L. \textit{Smart Drums: Comparing Audio and MIDI in Embedded Real-Time Drum Pattern Recognition}. Journal of the Audio Engineering Society, 2026 (in press).
    \item Silva, N. and Turchet, L. \textit{Music Information Retrieval as a Service: Latency Characterization of an Edge-to-Server Architecture for Near-Real-Time Audio Analysis}. IEEE International Symposium on the Internet of Sounds, 2026 (accepted).
    \item Silva, N., Wanderley, M. and Turchet, L. \textit{Melody and Motion: Integrating Guitar Gestures with Musical Patterns for Extended Control in Live Performance and Metaverse Applications}. IEEE International Symposium on the Internet of Sounds, Oct. 2025.
    \item Silva, N., Boem, A. and Turchet, L. \textit{Interactive IoMusT-Based Concerts: Real-Time Pattern Recognition and Audience Experience}. IEEE International Conference on Immersive and 3D Audio, Sep. 2025.
    \item Silva, N. and Turchet, L. \textit{User-Centered Evaluation of Smart Musical Instruments with Embedded Real-Time Pattern Detection}. IEEE International Conference on Immersive and 3D Audio, Sep. 2025.
    \item Silva, N. and Turchet, L. \textit{Real-Time Pattern Recognition of Symbolic Monophonic Music}. 19th International Audio Mostly Conference (ACM), Sep. 2024.
    \item Silva, N. and Turchet, L. \textit{A Structural Similarity Index Based Method to Detect Symbolic Monophonic Patterns in Real-Time}. 25th International Conference on Digital Audio Effects (DAFx), Sep. 2022.
    \item Silva, N., Fischione, C. and Turchet, L. \textit{Towards Real-Time Detection of Symbolic Musical Patterns: Probabilistic vs. Deterministic Methods}. 27th FRUCT Conference, 2020.
    \item Silva, N., Weeraddana, P.C. and Fischione, C. \textit{On Musical Onset Detection via the S-Transform}. 52nd Asilomar Conference on Signals, Systems, and Computers, Oct. 2018.
\end{itemize}

% ============================== DATASETS =============================
% [V2 keep in full] [V3 keep, or compress to the two most relevant]
% [V1/V4/V5/V6 DELETE and add this single line to Highlights instead:
%   \item \textbf{Open datasets:} Four open-access musical-pattern datasets
%   (DoMP, DoPP, DoDP, DoDP2) -- 9,800+ recordings from 70+ musicians,
%   published on Zenodo as benchmarks for real-time pattern detection.]
\section*{Datasets \footnotesize {\normalfont{(open access, Zenodo)}}}
\begin{itemize}[leftmargin=1.2em, itemsep=-0.25em]
    \item \textbf{DoMP} -- 4,392 monophonic MIDI recordings, 40 musicians; symbolic pattern-recognition benchmark (\href{https://doi.org/10.5281/zenodo.10818617}{\textit{10.5281/zenodo.10818617}})
    \item \textbf{DoPP} -- 2,276 polyphonic audio recordings (44.1 kHz/24-bit), 20 musicians; raw-audio RNN training and evaluation (\href{https://doi.org/10.5281/zenodo.14497998}{\textit{10.5281/zenodo.14497998}})
    \item \textbf{DoDP} -- 994 drum-pattern MIDI recordings, 10 drummers (\href{https://doi.org/10.5281/zenodo.14497974}{\textit{10.5281/zenodo.14497974}})
    \item \textbf{DoDP2} -- 2,177 drum-pattern MIDI recordings, 10 drummers, artist-level style templates (\href{https://doi.org/10.5281/zenodo.18395007}{\textit{10.5281/zenodo.18395007}})
\end{itemize}

% ==================== SCIENTIFIC COMMITTEE & PEER REVIEW =============
% [V2 keep] [V3 keep if the JD names peer review] [V1/V4/V5/V6 DELETE]
\section*{Scientific Committee and Peer Review}
\begin{itemize}[leftmargin=1.2em, itemsep=-0.25em]
    \item \textbf{Reviewer:} IEEE Access; Journal of the National Science Foundation, Sri Lanka; Colloquio di Informatica Musicale.
    \item \textbf{Program Committee:} IEEE International Symposium on the Internet of Sounds; International Conference on Immersive and 3D Audio.
\end{itemize}

% =========================== SELECTED PROJECTS =======================
% [V3 and V5 keep — this is where the domain evidence lands]
% [V1 keep only if Experience bullets did not already cover the JD's domain]
% [V2/V4/V6 DELETE]
% V3 draw from KG §8.1 (P1, P2, P3, P3b, P4); V5 from KG §8.2 (P14, P15, P16, P18, P19).
\section*{Selected Projects \footnotesize {\normalfont{(full portfolio: \href{https://nishal.xyz/research}{\textit{nishal.xyz/research}})}}}
\begin{itemize}[leftmargin=1.2em, itemsep=-0.25em]
    \item \textbf{<<KG §8 | project 1 name>>:} <<KG §8 | one sentence, lead with the metric>>
    \item \textbf{<<KG §8 | project 2 name>>:} <<KG §8 | ...>>
    \item \textbf{<<KG §8 | project 3 name>>:} <<KG §8 | ...>>
\end{itemize}

% =============== TALKS, DEMOS & INVITED PRESENTATIONS ================
% [V2 only, and only at 3 pages] Keep 5-8, most recent first (KG §9).
% \section*{Talks, Demos, and Invited Presentations}
% \begin{itemize}[leftmargin=1.2em, itemsep=-0.25em]
%     \item \textbf{Notte della Ricerca, MUSE Science Museum}, Trento, Italy -- Demo: smart musical instrument controlling peripheral devices. \hfill \textit{Sep 2025}
%     \item \textbf{Musical Metaverse (MUSMET) Kickoff}, University of Trento, Italy -- Mixed-reality demo for smart musical instruments. \hfill \textit{2025}
%     \item \textbf{University of Trento}, Italy -- Multisensory concert enabled by smart musical instruments and real-time pattern detection. \hfill \textit{Oct 2024}
%     \item \textbf{Audio Mostly Conference}, University of Milan, Italy -- Paper and system demo. \hfill \textit{Sep 2024}
%     \item \textbf{ICT Days}, University of Trento, Italy -- Real-time pattern detection demo. \hfill \textit{Sep 2024}
%     \item \textbf{McGill University}, Montreal, Canada -- Invited talk: real-time musical pattern detection. \hfill \textit{May 2024}
%     \item \textbf{ISMIR Conference}, Politecnico di Milano, Italy -- Smart musical instrument demo. \hfill \textit{Oct 2023}
%     \item \textbf{SoundMIT Synth Expo}, Turin, Italy -- Smart musical instrument demo. \hfill \textit{Oct 2023}
% \end{itemize}

% ======================== STUDENT SUPERVISION ========================
% [V2 only, or any JD naming mentorship/supervision] KG §4.1a.
% DATA PENDING - leave commented until the knowledge graph is filled in.
% Do not invent entries.
% \section*{Student Supervision}
% \begin{itemize}[leftmargin=1.2em, itemsep=-0.25em]
%     \item \textbf{<<KG §4.1a | name>>} (<<KG §4.1a | BSc/MSc/PhD/visiting/intern>>, <<KG §4.1a | institution>>) -- <<KG §4.1a | topic>>. \hfill \textit{<<KG §4.1a | year(s)>>}
%     \item \textbf{<<KG §4.1a | name>>} (<<KG §4.1a | level>>, <<KG §4.1a | institution>>) -- <<KG §4.1a | topic; note co-supervision or resulting publication>>. \hfill \textit{<<KG §4.1a | year(s)>>}
% \end{itemize}

% ===================== HIGHLIGHTS & RECOGNITION ======================
% [ALL VARIANTS] Last section. Keep 3-5 bullets; swap bullet 3 per variant.
\vspace{-1em}
\section*{Independent Research, Highlights, \& Recognition}
\begin{itemize}[leftmargin=1.2em, itemsep=-0.25em]
    \item \textbf{MIDI Innovation Awards:} Finalist (Sep 2023), Prototypes \& Non-Commercial Software category -- ``Hot Licks'' real-time pattern detection for smart instruments.
    \item \textbf{GMOA Recognition (Sri Lanka):} Computer-vision system for COVID-19 patient vitals monitoring deployed across hospital wards; enabled remote monitoring and reduced healthcare-worker exposure.
    \item <<VARIANT | V1 foundation-model music generation research | V2 the open-datasets line (see Datasets block) | V3 real-time ML prototypes for audio and sensor streams under constrained compute | V4 performing and recording guitarist, applying performer intuition to human-AI interaction in live music | V5 Promptly: on-demand printing product now deployed across four countries>>
    \item Volunteer guitarist for community music programs, performing regularly and collaborating with other musicians.
\end{itemize}

\end{document}
